\documentclass[12pt,a4paper]{article}
% ============================================================
%  Structural Persistence Series — Shared Preamble
% ============================================================
\usepackage{fontspec}
\setmainfont{Hiragino Mincho ProN}
\setsansfont{Hiragino Sans}
\setmonofont{Menlo}

\usepackage{iftex}
\ifXeTeX
  \XeTeXlinebreaklocale "ja"
  \XeTeXlinebreakskip=0pt plus 1pt minus 0.1pt
\fi

\usepackage[margin=22mm]{geometry}
\usepackage{amsmath,amssymb,amsthm}
\usepackage{xcolor}
\usepackage{graphicx}
\usepackage{hyperref}
\usepackage{booktabs}
\usepackage{fancyhdr}
% enumitem not available in this TeX installation

% --- Colours ------------------------------------------------
\definecolor{PaperAccent}{HTML}{1F4E79}
\definecolor{PaperGray}{HTML}{51606F}

\hypersetup{
  colorlinks=true,
  linkcolor=PaperAccent,
  urlcolor=PaperAccent,
  citecolor=PaperAccent,
  pdflang=ja
}
\urlstyle{same}

% --- Typography ---------------------------------------------
\setlength{\parindent}{1.5em}
\setlength{\parskip}{0pt}
\setlength{\emergencystretch}{3em}
\linespread{1.08}
\setlength{\abovedisplayskip}{8pt plus 2pt minus 4pt}
\setlength{\belowdisplayskip}{8pt plus 2pt minus 4pt}
\setlength{\abovedisplayshortskip}{6pt plus 2pt minus 2pt}
\setlength{\belowdisplayshortskip}{6pt plus 2pt minus 2pt}

% --- Section label names ------------------------------------
\renewcommand{\abstractname}{要旨}

% --- Theorem-like environments ------------------------------
\theoremstyle{plain}
\newtheorem{proposition}{命題}[section]
\newtheorem{lemma}[proposition]{補題}
\newtheorem{corollary}[proposition]{系}

\theoremstyle{definition}
\newtheorem{definition}{定義}[section]
\newtheorem{assumption}{条件}[section]

\theoremstyle{remark}
\newtheorem*{remark}{注}

% --- Running header -----------------------------------------
\pagestyle{fancy}
\fancyhf{}
\renewcommand{\headrulewidth}{0.4pt}
\fancyhead[L]{\small\textcolor{PaperGray}{\nouppercase{\leftmark}}}
\fancyhead[R]{\small\textcolor{PaperGray}{\thepage}}
\fancypagestyle{plain}{%
  \fancyhf{}
  \renewcommand{\headrulewidth}{0pt}
  \fancyfoot[C]{\small\textcolor{PaperGray}{\thepage}}
}

% --- Title block macro --------------------------------------
\newcommand{\PaperTitleBlock}[3]{%
  % #1 = title, #2 = subtitle, #3 = date
  \thispagestyle{plain}%
  \begin{center}
  {\LARGE\bfseries #1\par}
  \vspace{0.4em}
  {\normalsize #2\par}
  \vspace{1.0em}
  {\large Akihito Sunagawa\par}
  \vspace{0.3em}
  {\normalsize #3\par}
  \end{center}
  \vspace{1.6em}
}

% Legacy 2-arg version (backward compat)
\newcommand{\PaperTitleBlockSimple}[2]{%
  \PaperTitleBlock{#1}{}{#2}
}


\begin{document}

\PaperTitleBlock
  {構造持続理論の統合版}
  {— 推論劣化・継続学習忘却・持続知能アーキテクチャ —}
  {2026年4月12日}

\begin{abstract}
本稿は、構造持続理論に関する既存の分冊を、一つの導線で読めるように再構成した統合版である。中心にある主張は、長期対話における推論劣化と、継続学習における破滅的忘却とが、いずれも「構造を維持できる状態の領域が、未整理の矛盾や前提更新によって縮小していく過程」として記述できる、という点にある。

理論の最小核では、構造維持可能な状態集合の逐次縮小を対数比で測ることで、残存可能性が指数型で表されることを与える。したがって、指数型は追加仮定ではなく、この表現の帰結である。条件つき導出では、どこまでが定義の帰結であり、どこからが独立性や弱依存といった追加条件に依存するのかを明示する。

推論側では、長期対話の劣化を文脈長そのものではなく、未整理の矛盾蓄積として捉える。継続学習側では、前提更新に依存する知識の連鎖更新が失敗することを、構造的忘却として記述する。さらに、これらの含意として、外部代謝、多層記憶、ロールバック可能性を備えた持続知能アーキテクチャの方向を述べる。本稿は、理論、実験、実装含意を一つの見取り図として提示するための入口であり、詳細な導出と実験条件は各分冊に委ねる。
\end{abstract}

\section{はじめに}\label{ux306fux3058ux3081ux306b}

大規模言語モデルにおける長期的一貫性の問題は、ふつう別々の話題として扱われる。一つは、長い対話のなかで矛盾や上書きが蓄積し、推論が崩れていく問題である。もう一つは、新しい知識を逐次学習させると古い知識が壊れてしまう、いわゆる破滅的忘却である。

本稿群は、LLM を一種の構造維持系として見たとき、長期対話での推論劣化と継続学習での忘却は同じ骨格の問題として現れているのではないかと考え、その見方がどこまで理論的・実験的に支えられるかを確かめるために書かれた。

本稿の立場では、これらは無関係な二問題ではない。どちらも、ある構造を維持できる状態の領域が、制約の蓄積によって縮小していくという、より小さな骨格の上で読むことができる。推論時には、構造とは「与えられた前提や規則のもとで、論理一貫性を保って最後まで到達できること」である。継続学習時には、構造とは「更新された前提のもとで、依存知識が整合的に再編されていること」である。

ここで問題にしているのは、基体そのものの消滅ではない。問題にしているのは、ある構造としての同一性や、その構造が担っていた機能が持続できるかどうかである。したがって、劣化や崩壊として観測されるものは、多くの場合、無への消失ではなく、別の構造への遷移、構造的再編、相転移として理解される。

本稿の目的は、この共通骨格を最小理論、条件つき導出、推論実験、継続学習実験、実装含意の順に再配置し、全体像を一つの物語として読めるようにすることにある。新しい定理や新しい実験を追加することが目的ではない。目的は、既存の分冊がそれぞれ担っていた役割を、一つの統一的主張のもとに並べ直すことである。

\section{最小理論}\label{ux6700ux5c0fux7406ux8ad6}

系に対し、その構造を維持できる状態の集合を \[
V^{(0)}
\]

とする。制約追加モデル \[
V^{(i)} = V^{(i-1)} ∩ C_{i}
\]

のもとでは、 \[
V^{(0)} \supseteq V^{(1)} \supseteq \cdots \supseteq V^{(n)}
\]

という縮小列は命題として従う。より一般には、学習・修復・外部支援・ロールバックのような再拡大過程もありうるため、本稿群ではこの縮小性を収縮モードに対する適用条件として読む。理論の主対象は、残存比率 r ∈ (0, 1{]} 上の損失尺度 f : (0, 1{]} → {[}0, ∞) であり、集合対 \(A \supseteq B\)に対する損失 D(A, \$B) = \(f(m(B)/m(A)\))\$はその表現として読まれる。損失尺度に対する自然な公理系（比率依存性、正規化、加法性、連続性）を置くと、f は対数比として一意に強制される（分冊 1 §3 の対数比の一意性定理）。単位規約 \(k = 1\)のもとで、各段階の損失は \[
l_{i} = -\ln( m(V^{(i)}) / m(V^{(i-1)}) )
\]

となる。累積損失 \[
L_{n} = \sum l_{i}
\]

に対して \[
m(V^{(n)}) = m(V^{(0)}) e^{-L_{n}}
\]

が恒等式として成り立つ。

したがって、指数型は経験的に当てはめた関数形ではなく、損失測度の公理系から二段階で導かれる必然的な形式である。対数比の一意性は、Hartley \((1928)\) / Shannon \((1948)\) による情報量の特徴づけと同じ論法（Cauchy の関数方程式の連続解）を、確率測度ではなく構造維持可能性の測度 m に適用したものにあたる。加法性の要請（B3）は、独立サブシステムでの損失が足し合わさるという示量性を自然な動機づけとして持ち、統計力学におけるエントロピーの示量性と同じ骨格を共有する。各ドメインは f 全域ではなく、そのドメインで実現される比率の部分集合を観測するものと位置づけられるため、離散的な測度空間で f の一部が未実現のまま残ることは理論の限界ではなく、f を各ドメインが部分的にサンプリングすることの反映として読む。この最小核に、有効維持資源 M を別に導入すれば、構造持続ポテンシャル \[
S = M e^{-L}
\]

を得る。ここで重要なのは、M が資源側、L が構造維持可能性の縮小側を表し、この二つを切り分けられることである。資源が残っていても構造が壊れる場合があり、逆に構造を保てる余地がなお残っていても資源不足で維持できない場合がある。

この最小理論が与える予測は、制約の個数そのものよりも、どの制約がどれだけ構造維持可能状態を削るかという「質」の差に関する順序予測である。以後の各分冊は、この最小理論をどのように具体化できるかを別々の文脈で問うている。

この点は、同じ基体が別の構造として存続しうる場合にとくに重要である。たとえば同じ分子組成が液体・固体・気体として異なる相に現れるとき、何を維持すべき構造とみなすかによって、観測されるのは消滅ではなく相の遷移になる。本稿群が扱うのは、そのような「何として持続しているのか」を問うための最小骨格である。

\section{条件つき導出と理論の射程}\label{ux6761ux4ef6ux3064ux304dux5c0eux51faux3068ux7406ux8ad6ux306eux5c04ux7a0b}

最小理論の指数式そのものは、状態集合の縮小率を対数比で定義する限り恒等式であり、独立性を要しない。しかし、各段階損失の生成過程を確率的にモデル化するときには、どこからが追加条件に依存するかを分けておく必要がある。

このため統合版では、三つの条件を区別する。第一に、制約が縮小方向に働くこと。これは制約追加モデルでは命題として従うが、より一般にはなお適用条件として読むべきものである。第二に、損失を縮小率の対数比で測ること。ここまでで恒等式としての指数表現は成立する。第三に、段階損失の生成過程に独立性または弱依存の制御を入れること。これは恒等式のためではなく、参照モデルに対して指数型の安定性を論じるための条件である。

この区別は理論の言い過ぎを避けるために重要である。本稿群が直接に言っているのは、「あらゆる実系で構造持続が厳密に指数減衰する」という強い普遍命題ではない。より正確には、「最小形式の指数型は定義の帰結であり、実系への適用では、どの測度を選び、どの制約がどの縮小率を持つか、またその生成過程がどこまで安定に制御されるかが問題になる」ということである。

したがって、この理論の強さは、指数型そのものではなく、個別ドメインにおいて基準モデルを上回る追加の予測力を持てるかどうかにかかる。推論時の矛盾蓄積、継続学習時の前提更新、SAT における探索難度などは、その具体化候補として並んでいる。

\section{推論時の構造劣化}\label{ux63a8ux8ad6ux6642ux306eux69cbux9020ux52a3ux5316}

推論側の分冊では、長い対話で起きる劣化を、単なる文脈長の問題としてではなく、未整理の矛盾や更新が蓄積することで、有効な推論経路の集合が縮小していく過程として記述した。ここでいう有効経路とは、与えられた前提や規則に反せず、自己矛盾なく最後まで到達できる粗い推論段階列の集合である。

この設定では、構造的な矛盾、指示衝突、時間的な不整合、参照元どうしの衝突、未整理の上書きなどが制約として働き、 \[
A^{(0)} \supseteq A^{(1)} \supseteq \cdots \supseteq A^{(n)}
\]

という有効経路集合の縮小を生む。したがって、経路集合の残存量もまた最小理論と同じ指数形式で表される。

実験では、矛盾する更新を外部で検出し、「旧値 -\textgreater{} 新値」の時間ラベルつき対として整理・保持する条件を ON、同じ矛盾を注入しつつその整理を行わない条件を OFF、矛盾注入なしを NC とした。観測された主結果は、ON が OFF を安定して上回り、OFF が NC を下回るという方向である。ここから得られる最小の経験的主張は、長期対話の劣化は単に長さだけでは説明しにくく、未整理の矛盾が有力な制約源になっている、という点である。

この分冊の役割は、理論が推論全般を説明すると言うことではなく、少なくとも論理一貫性維持という限定された観測量に関して、構造持続の枠組みが基準モデルに対する追加予測力を持つことを示す最初の実験を与えることであった。

\section{継続学習時の構造的忘却}\label{ux7d99ux7d9aux5b66ux7fd2ux6642ux306eux69cbux9020ux7684ux5fd8ux5374}

継続学習側の分冊では、忘却を単に「古い事実を思い出せないこと」としてではなく、前提更新に依存する知識の連鎖更新が失敗し、内部状態が新旧の知識を混在させた不整合になることとして捉えた。これを構造的忘却と呼ぶ。

たとえば「会社 X の本社は東京である」が「大阪である」に更新されたなら、最寄駅、通勤路線、配送区域など、その前提に依存する知識も同時に組み替えられなければならない。上流の前提だけが更新され、下流が古いまま残るなら、それは単なる旧知識喪失ではなく、構造的再編の失敗である。

この問題に対し、分冊では LoRA ベースの逐次学習を用い、過去再提示の基準条件 E-lite、前提依存構造を追跡する F-v2c、現在知識と過去知識を別アダプタに分ける F-multi を比較した。結果は、全モデル・全条件で最初の前提更新以降に旧知識保持が急減し、LoRA が蓄積というより上書きに近く振る舞うことを示した。F-v2c は依存整合性を改善するが旧知識保持忠実度は低く、F-multi は非ゼロの旧知識保持を与えるものの、高忠実度の長期保持にはなお届かない。

この分冊の役割は、推論側で扱った「未整理の矛盾が経路を削る」という話を、学習側では「前提更新が依存構造を壊す」という別の相として示すことである。共通点は、どちらも構造を維持できる状態の領域が縮小していくことにある。

\section{実装含意}\label{ux5b9fux88c5ux542bux610f}

理論と実験が指しているのは、単にモデルを大きくすればよい、あるいは訓練レシピを改善すればよい、というだけの方向ではない。少なくとも本稿群が扱った現象は、未整理の矛盾を抱えたまま長期に運用する stateless な連想機械の限界を示唆している。

そこから出てくる設計含意は、外部代謝、多層記憶、ロールバック可能性である。外部代謝とは、矛盾する更新を「古い状態」と「新しい状態」の関係として外で整理し、未整理の上書きのまま会話や学習に残さないことである。多層記憶とは、作業記憶、活性ルール、休眠知識、事実アーカイブを分け、同じ形式ですべてを抱え込まないことである。ロールバック可能性とは、代謝や更新が失敗したときに、局所的に復元できることである。

この方向は、長期的な一貫性を持つ知能システム、すなわち内部モデルを育てながら持続するパートナー型知能を考えるときに自然である。記憶を増やすことそのものよりも、矛盾をどう解消し、何を休眠させ、いつ再活性化し、どこまで安全に戻せるかが中心になるからである。

\section{限界と今後の課題}\label{ux9650ux754cux3068ux4ecaux5f8cux306eux8ab2ux984c}

本稿群は、構造持続理論のすべてを証明したものではない。最小形式の指数型は定義の帰結であり、その意味で骨格は明確だが、各ドメインで何を構造維持可能状態とみなすか、どの測度で比較するか、どの代理指標が最もよいかはなお個別課題として残る。

推論実験について言えば、現時点で強く言えるのは \(ON > OFF\)と \(OFF < NC\)の方向であり、ON が常に NC を上回ることまで一般化するのは早い。継続学習実験について言えば、LoRA ベース逐次更新という一つの更新様式の限界を示したのであって、あらゆる訓練手法を尽くしたわけではない。また、構造的忘却の記述はなお操作的であり、内部表現そのものを直接観測しているわけではない。

それでも、本稿群の価値は、推論劣化と継続学習忘却を別々の現象として終わらせず、どちらも構造維持可能性の縮小として読む一つの言語を与えた点にある。今後の課題は、この言語をより強い Route A の定理ルートと、より厚い Route C の実証ルートの両方で鍛えることである。

\section{結論}\label{ux7d50ux8ad6}

本稿は、構造持続理論に関する既存の分冊を統合し、長期対話における推論劣化と、継続学習における破滅的忘却とを、一つの最小骨格の上に再配置した。最小理論では、構造を維持できる状態集合の縮小から指数型が現れることを与えた。条件つき導出では、その指数型がどこまで定義の帰結で、どこからが追加条件に依存するのかを分けた。推論実験では、未整理の矛盾が論理一貫性を保てる経路を削ることを示した。継続学習実験では、前提更新が依存知識の再編を必要とし、それが失敗すると構造的忘却が起きることを示した。

したがって、本稿群の中心主張は次のように要約できる。長期的な知能システムの劣化を理解する鍵は、単なる容量不足や長さそれ自体ではなく、構造を維持できる状態の領域がどのような制約によって削られていくかを見ることにある。そこから自然に導かれる設計方向は、矛盾を外部で代謝し、内部モデルを壊れにくく保つ持続知能アーキテクチャである。

ここで失われるのは、存在一般ではなく、「その構造として持続すること」である。この整理を置いておくことで、崩壊、忘却、相転移、再編といった異なる観測現象を、同じ骨格の別相として読むことができる。

\section{詳細への案内}\label{ux8a73ux7d30ux3078ux306eux6848ux5185}

本稿は入口であり、詳細は各分冊に委ねる。

\begin{itemize}
\tightlist
\item
  最小理論: 1\_構造持続の最小形式.md
\item
  条件つき導出: 2\_構造持続の条件つき導出.md
\item
  推論実験: 3\_構造持続と推論性能の劣化.md
\item
  継続学習実験: 4\_構造持続と継続学習における破滅的忘却.md
\item
  実装含意: 補論\_持続的支援知能の設計原理.md
\item
  追加補論: 補論\_構造持続写像の標準手順.md, 補論\_計算コストの構造的予測.md
\end{itemize}
\end{document}
